\section{Doubly Robust Adjustment for Local-Law Losses}
\label{app:doubly-robust-local-law-adjustment}

The local laws are stated as oracle conditions: a state should preserve
the value that the task oracle \(\fstar\) assigns to the object it
represents. In the cleanest version, every local check can be evaluated
with \(\fstar\). In the teacher-first regime this is not what we see.
We usually have document-level labels at the root, use those labels to
fit or select a learned query \(\widehat f\), and then evaluate most
local-law violations through \(\widehat f\). Direct node-level oracle
feedback, when available, arrives only on a sampled subset of nodes.

This creates a measurement problem. The proxy local-law loss measured
through \(\widehat f\) is useful because it is available everywhere, but
it need not equal the local-law loss that would have been measured
through \(\fstar\). The correction below is the standard DSL
design-adjusted outcome applied to the local-law loss itself. It does
not estimate the full oracle function \(\fstar\). It estimates the
oracle-measured local-law objective that would be used if every node
had oracle feedback. To keep the notation readable, we use descriptive
names: \(L_\star\) for the oracle law loss, \(L_{\mathrm{proxy}}\) for
the all-proxy law loss, and \(L_{\mathrm{corrected}}\) for the
corrected law loss actually used by the objective. We write
\(L_{\mathrm{gap}}\) for the proxy-oracle gap loss, the loss that
measures how far the learned query \(\widehat f\) is from \(\fstar\) on
the available root-level supervision.

\subsection{The Estimand}
\label{app:drlaw-estimand}

Fix a realized tree \(T\). For each node \(v\), let \(d(v)\) be its
depth, with \(d(v)=0\) at the root. Suppress the C1/C2/C3 channel index
for notational clarity; the same construction applies separately to
each law channel, or to any chosen scalar mixture of them. Let
\[
  \ell_v^\star
  =
  \ell_{\mathrm{law}}(v;\fstar,g)
\]
be the local-law loss at node \(v\) measured with the oracle, and let
\[
  \widehat\ell_v
  =
  \ell_{\mathrm{law}}(v;\widehat f,g)
\]
be the corresponding proxy loss measured with the learned query. The
target local-law functional is the discounted oracle law loss
\begin{equation}
\label{eq:drlaw-oracle-objective}
  L_\star(T)
  =
  \sum_{v\in V(T)}
  \gamma^{d(v)}\,\ell_v^\star .
\end{equation}
This is the quantity the local-law part of the objective would use if
all node-level oracle labels were available.

The root supervision term is separate. It trains or evaluates the final
answer at the root using the available document-level label. The local
law term \(L_\star(T)\) instead asks whether the
realized intermediate states behave like oracle-valid summaries.

The learned query \(\widehat f_\theta\) also has its own proxy-oracle
gap loss, which is the clearer name for the term sometimes called
calibration in older notes:
\[
  L_{\mathrm{gap}}(\theta)
  =
  \E_{x\in\mathcal D_{\mathrm{root}}}
  \ell_{\mathrm{gap}}
  \bigl(\widehat f_\theta(x),\fstar(x)\bigr).
\]
For a metric-valued oracle one may take
\(\ell_{\mathrm{gap}}(\widehat f_\theta(x),\fstar(x))
=d_Y(\widehat f_\theta(x),\fstar(x))\), or the corresponding squared
or task-specific root-label loss. This term measures how well
\(\widehat f\) matches the oracle used to define local laws; it is not
the root-output loss of the tree itself.

\subsection{Node-Level DSL Adjustment}
\label{app:drlaw-node-adjustment}

Let \(R_v\in\{0,1\}\) indicate whether node \(v\) receives direct
oracle feedback, and let \(\pi_v>0\) be the propensity used for the
inverse-propensity correction. The corrected node loss is
\begin{equation}
\label{eq:drlaw-node-adjustment}
  \ell_v^{\mathrm{corrected}}
  =
  \widehat\ell_v
  +
  \frac{R_v}{\pi_v}
  \left(\ell_v^\star-\widehat\ell_v\right).
\end{equation}
Equivalently,
\[
  \ell_v^{\mathrm{corrected}}
  =
  \widehat\ell_v
  -
  \frac{R_v}{\pi_v}
  \left(\widehat\ell_v-\ell_v^\star\right),
\]
which is exactly the DSL design-adjusted outcome with
prediction \(\widehat\ell_v\) and observed oracle outcome
\(\ell_v^\star\). If \(R_v=0\), the corrected loss is the proxy loss.
If \(R_v=1\) and \(\pi_v=1\), the corrected loss is the oracle loss.

It is often easiest to read this through the proxy-oracle local-law
bias
\[
  b_v
  =
  \widehat\ell_v-\ell_v^\star .
\]
Then the same corrected node loss is
\begin{equation}
\label{eq:drlaw-node-bias-form}
  \ell_v^{\mathrm{corrected}}
  =
  \ell_v^\star
  +
  \left(1-\frac{R_v}{\pi_v}\right)b_v .
\end{equation}
Thus the correction does not hide the bias: it leaves exactly the
residual multiplier \(1-R_v/\pi_v\) times the proxy-oracle law error.

Condition on the realized tree, the fitted \(g\), the fitted
\(\widehat f\), and all logged design information except the Bernoulli
draw \(R_v\). Let
\[
  \pi_v^\star
  =
  \Prob(R_v=1\mid \mathcal F)
\]
be the true marginal inclusion probability under this conditioning
information. The DSL misspecified-propensity identity gives the exact
residual:
\begin{align}
  \E\!\left[
    \ell_v^{\mathrm{corrected}}-\ell_v^\star
    \mid \mathcal F
  \right]
  &=
  \E\!\left[
    \left(1-\frac{R_v}{\pi_v}\right)
    \left(\widehat\ell_v-\ell_v^\star\right)
    \mid \mathcal F
  \right] \notag\\
  &=
  \left(1-\frac{\pi_v^\star}{\pi_v}\right)
  \left(\widehat\ell_v-\ell_v^\star\right).
\label{eq:drlaw-misspecified-residual}
\end{align}
The first equality is algebra from~\eqref{eq:drlaw-node-adjustment};
the second is the design-adjusted DSL expectation identity. No
independence across nodes is needed for this scalar statement. For the
tree aggregate, linearity of expectation applies node by node.

\subsection{Doubly Robust Reading}
\label{app:drlaw-doubly-robust}

Equation~\eqref{eq:drlaw-misspecified-residual} is the doubly robust
statement in this setting. The residual vanishes under either of two
routes:
\[
  \pi_v=\pi_v^\star
  \qquad\text{or}\qquad
  \widehat\ell_v=\ell_v^\star .
\]
The first route is correct design logging: the propensity used in the
correction is the actual node inclusion probability. Then the oracle
residual is removed regardless of the quality of \(\widehat f\). The
second route is exact local-law prediction: the proxy loss already
equals the oracle loss, so there is no proxy error left for a wrong
propensity to amplify.

When both routes fail, the remaining bias is explicit:
\[
  \left(1-\frac{\pi_v^\star}{\pi_v}\right)
  \left(\widehat\ell_v-\ell_v^\star\right).
\]
This is why the propensity residual is kept visible in the local-law
certificate. Sampling that depends on proxy violations, depth, node
type, retry behavior, or human triage is harmless only if the logged
propensity represents that realized sampling law. If the logged
propensity ignores part of the sampling mechanism, the residual above
is the price paid, scaled by the proxy-oracle local-law error.

\subsection{Discounted Tree Aggregation}
\label{app:drlaw-tree-aggregation}

The corrected tree-level law loss is
\begin{equation}
\label{eq:drlaw-adjusted-tree-loss}
  L_{\mathrm{corrected}}(T)
  =
  \sum_{v\in V(T)}
  \gamma^{d(v)}\,\ell_v^{\mathrm{corrected}} .
\end{equation}
Substituting~\eqref{eq:drlaw-node-adjustment} gives the simple
one-line adjustment:
\[
  L_{\mathrm{corrected}}(T)
  =
  \sum_{v\in V(T)}
  \gamma^{d(v)}
  \left[
    \widehat\ell_v
    +
    \frac{R_v}{\pi_v}
    \left(\ell_v^\star-\widehat\ell_v\right)
  \right].
\]
The endpoints match the intended cases. If there are no node-level
oracle observations, \(R_v=0\) for every node, and
\[
  L_{\mathrm{corrected}}(T)
  =
  L_{\mathrm{proxy}}(T)
  =
  \sum_{v\in V(T)}
  \gamma^{d(v)}\,\widehat\ell_v ,
\]
the proxy-only discounted local-law objective. If every node is
observed with \(\pi_v=1\), then
\[
  L_{\mathrm{corrected}}(T)
  =
  \sum_{v\in V(T)}
  \gamma^{d(v)}\,\ell_v^\star
  =
  L_\star(T).
\]
Under correct logged propensities, the conditional expectation of
the corrected aggregate equals the oracle aggregate:
\[
  \E\!\left[
    L_{\mathrm{corrected}}(T)
    -
    L_\star(T)
    \mid \mathcal F
  \right]
  =
  0.
\]
Under misspecified propensities, the aggregate residual is the
discounted sum of the node residuals in
\eqref{eq:drlaw-misspecified-residual}.

Using the bias notation \(b_v=\widehat\ell_v-\ell_v^\star\), the same
tree loss can be written as
\begin{equation}
\label{eq:drlaw-tree-bias-form}
  L_{\mathrm{corrected}}(T)
  =
  L_\star(T)
  +
  \sum_{v\in V(T)}
  \gamma^{d(v)}
  \left(1-\frac{R_v}{\pi_v}\right)b_v .
\end{equation}
Taking conditional expectations replaces \(R_v\) by
\(\pi_v^\star\), so the expected residual bias is
\[
  \E\!\left[
    L_{\mathrm{corrected}}(T)-L_\star(T)
    \mid\mathcal F
  \right]
  =
  \sum_{v\in V(T)}
  \gamma^{d(v)}
  \left(1-\frac{\pi_v^\star}{\pi_v}\right)b_v .
\]
This vanishes when propensities are logged correctly, and it also
vanishes when the proxy law loss is exact, \(b_v=0\), even if the
propensity is misspecified.

\subsection{From Adjustment to the Single-\texorpdfstring{\(\lambda\)}{lambda} Objective}
\label{app:drlaw-single-lambda}

The corrected law loss is the local-law channel used by the objective.
There is no additional agreement gate and no additive gap penalty in
the optimized loss. Restoring the candidate parameter \(\theta\) and
the supervised training channels, the theorem-facing objective is the
two-channel root/local objective
\begin{equation}
\label{eq:drlaw-final-objective}
  J(\theta;T)
  =
  \underbrace{
  (1-\lambda)
  L_{\mathrm{root}}(\theta)
  }_{\text{root supervision}}
  +
  \underbrace{
    \lambda L_{\mathrm{corrected}}(\theta,T)
  }_{\text{corrected local laws}} .
\end{equation}
In lambda mode, the user supplies \(\lambda\in[0,1]\) and an active law
set. The root channel receives share \(1-\lambda\), and the active
local laws split \(\lambda\) equally. If \(\lambda=0\), the objective is
root-only. In explicit-weight mode, the user instead supplies root and
per-law weights \(w_{\mathrm{root}}, w_1,\ldots,w_k\), supplies no
\(\lambda\), and we normalize
\[
  W = w_{\mathrm{root}}+\sum_{i=1}^k w_i,
  \qquad
  J(\theta;T)
  =
  \frac{w_{\mathrm{root}}}{W}L_{\mathrm{root}}(\theta)
  +
  \sum_{i=1}^k \frac{w_i}{W}L_i^{\mathrm{corrected}}(\theta,T).
\]
The implied local-law share is reported as
\(\lambda=\sum_i w_i/W\), but it is not accepted as a second input in
explicit-weight mode. A configuration that supplies both \(\lambda\)
and explicit root/law weights is therefore ill-posed and rejected.

The proxy-oracle gap, sampling uncertainty, oracle-transfer slack, and
propensity residual remain audit diagnostics for
\(L_{\mathrm{corrected}}\). They explain the quality of the corrected
local-law estimate, but they do not appear as additive terms in \(J\)
and they do not shrink \(\lambda\).

Substituting the bias form~\eqref{eq:drlaw-tree-bias-form} gives the
equivalent bias-explicit objective
\begin{align}
\label{eq:drlaw-final-objective-bias-form}
  J(\theta;T)
  &=
  \underbrace{
    (1-\lambda)
    L_{\mathrm{root}}(\theta)
  }_{\text{root supervision}}
  \notag\\
  &\quad+
  \lambda
  \Bigg[
    \underbrace{L_\star(\theta,T)}_{\text{oracle law target}}
    \notag\\
  &\qquad\qquad+
    \underbrace{
      \sum_{v\in V(T)}
      \gamma^{d(v)}
      \left(1-\frac{R_v}{\pi_v}\right)b_v(\theta)
    }_{\text{residual law bias}}
  \Bigg].
\end{align}
Equivalently, the last channel is
\[
  \underbrace{
    \lambda L_\star(\theta,T)
  }_{\text{oracle local-law pressure}}
  +
  \underbrace{
    \lambda
    \sum_{v\in V(T)}
    \gamma^{d(v)}
    \left(1-\frac{R_v}{\pi_v}\right)b_v(\theta)
  }_{\text{bias left after correction}} .
\]
This is the same objective as~\eqref{eq:drlaw-final-objective}. The
first term learns from root labels, and the bracketed local-law term is
the oracle local-law target plus whatever residual bias remains after
node-oracle correction. The proxy-oracle gap affects diagnostics for
the corrected channel, not the displayed objective weight.

\subsection{Interpretation}
\label{app:drlaw-interpretation}

The distinction between \(\fstar\) and the local-law estimand is
important. The corrected law loss above targets
\(L_\star(T)\), the oracle-measured local-law objective on a realized
tree. It does not recover the entire function \(\fstar\). The role of
\(\widehat f\) is as a nuisance/proxy: it provides local-law losses
everywhere and reduces variance, while direct node-oracle observations
correct its residual bias when sampled.

If the oracle is exactly mergeable within the chosen representation
class, correct sampling makes the local-law channel behave as if the
oracle local-law objective were observed. If the oracle is not
mergeable in the chosen class, nonzero local-law pressure does not
recover \(\fstar\) exactly; it projects the learned tree toward the
best local-law-compatible approximation available in the class. Oracle
labels affect this pressure through \(L_{\mathrm{corrected}}\), not
through a separate weight-shrinkage mechanism.
